\documentclass[11pt]{article}

\usepackage{amsmath,amssymb,amsthm,mathtools}
\usepackage[margin=1in]{geometry}
\usepackage{microtype}
\usepackage[hidelinks]{hyperref}

\newtheorem{theorem}{Theorem}[section]
\newtheorem{lemma}[theorem]{Lemma}
\newtheorem{corollary}[theorem]{Corollary}
\newtheorem{remark}[theorem]{Remark}

\newcommand{\dd}{\,\mathrm d}
\newcommand{\one}{\mathbf 1}

\title{Atlas 148: Sharp Paw Bias and a Two-Dimensional Hilbert Projection}
\author{}
\date{}

\begin{document}

\maketitle

\begin{abstract}
We prove the chromatic-polynomial lower bound for NetworkX Atlas graph 148
on the full interval $p\in[1/2,1]$ and for every graphon.  On
$[1/2,3/5]$, a Hilbert projection reduces its density to the square of a
paw density.  Symmetrizing the paw leaf, biasing by $d^{1/3}$, and applying
the sharp triangle-density theorem gives the required strengthened paw
bound.  On $[3/5,1]$, a two-dimensional projection retains the covariance
term that is lost by the false product comparison with two triangles and a
$4$-cycle.  The remaining scalar supporting line is certified exactly on
four rational Bernstein boxes.
\end{abstract}


%%%%%%%%%%%%%%%%%%%%%%%%%%%%%%%%%%%%%%%%%%%%%%%%%%%%%%%%%%%%%%%%%%%%%%%%%%%%%%%
\section{The graph, kernels, and target}
%%%%%%%%%%%%%%%%%%%%%%%%%%%%%%%%%%%%%%%%%%%%%%%%%%%%%%%%%%%%%%%%%%%%%%%%%%%%%%%

Let $(\Omega,\mathcal A,\mu)$ be a probability space and let
$W\colon\Omega^2\to[0,1]$ be symmetric and jointly measurable.  Put
\[
 p=\int_{\Omega^2}W(x,y)\dd\mu^2,
 \qquad
 d(x)=\int_\Omega W(x,y)\dd\mu(y),
\]
and define the common-neighbour and triangle-edge kernels
\[
 S(x,y)=\int_\Omega W(x,z)W(y,z)\dd\mu(z),
 \qquad K(x,y)=W(x,y)S(x,y).
\]

NetworkX Atlas 148, with graph6 code \texttt{EyUG}, has edge list
\[
 01,02,05,12,13,14,34,45.
\]
It consists of the triangles $012$ and $134$, which share vertex $1$,
together with the two-edge path $0$--$5$--$4$ between the indicated outer
vertices.  The triangle gives chromatic number at least three, while the
color classes $\{0,3\},\{1,5\},\{2,4\}$ give a proper three-colouring.
Thus its chromatic number is three and the admissible interval is
$p\in[1/2,1]$.

With $x=1$, $a=0$, and $b=4$, integrating vertices $2,3,5$ gives
the exact identity
\begin{equation}
 t(F_{148},W)
 =\int_{\Omega^3}K(x,a)K(x,b)S(a,b)\dd\mu(x)\dd\mu(a)\dd\mu(b).
 \label{eq:density-kernels}
\end{equation}
To count proper $r$-colourings, first colour $x$, then the two outer vertices
$a,b$, then the remaining vertex of each triangle, and finally the midpoint
of the outer path.  If $a,b$ have the same colour, the last vertex has
$r-1$ choices; if they have distinct colours, it has $r-2$ choices.  Hence
the count is
$r(r-2)^2\bigl((r-1)^2+(r-1)(r-2)^2\bigr)$, and the chromatic polynomial
and target are
\begin{align}
 \chi_{F_{148}}(r)
 &=r(r-1)(r-2)^2(r^2-3r+3),
 \label{eq:chromatic}\\
 \Phi_{F_{148}}(p)
 &=p(2p-1)^2(3p^2-3p+1).
 \label{eq:target}
\end{align}
Write
\[
 q=1-p,\qquad c=2p-1,\qquad f=3p^2-3p+1=p^3+q^3.
\]


%%%%%%%%%%%%%%%%%%%%%%%%%%%%%%%%%%%%%%%%%%%%%%%%%%%%%%%%%%%%%%%%%%%%%%%%%%%%%%%
\section{The normalized sharp triangle profile}
%%%%%%%%%%%%%%%%%%%%%%%%%%%%%%%%%%%%%%%%%%%%%%%%%%%%%%%%%%%%%%%%%%%%%%%%%%%%%%%

We use the graphon form of the clique density theorem for triangles.  Define
$g\colon[0,1]\to[0,1]$ as follows.  Choose an integer $k\geq1$ and
$u\in[1/(k+1),1/k]$ such that
\begin{equation}
 s=1-ku^2-(1-ku)^2,
 \label{eq:triangle-parameter}
\end{equation}
put
\begin{equation}
 g(s)=1-3\bigl(ku^2+(1-ku)^2\bigr)
       +2\bigl(ku^3+(1-ku)^3\bigr),
 \label{eq:triangle-profile}
\end{equation}
and set $g(1)=1$.  The definitions agree at the common endpoints.

\begin{theorem}[Sharp triangle density]
\label{thm:clique-density}
Every graphon $V$ of edge density $s$ satisfies
\[
 t(K_3,V)\geq g(s).
\]
\end{theorem}

This is the $K_3$ instance of the clique density theorem.  Its usual finite
asymptotic and standard-graphon forms imply the stated arbitrary-probability-
space form by sampling independent points from the given space and then
independent Bernoulli edges.  Edge and triangle densities converge in
probability to their graphon values, so no atomlessness assumption is needed.

The feature needed here is an elementary consequence of the explicit sharp
profile.

\begin{lemma}[Normalized triangle monotonicity]
\label{lem:normalized-triangle}
The function $g(s)/s^{3/2}$ is nondecreasing on $(0,1]$.
\end{lemma}

\begin{proof}
On the branch \eqref{eq:triangle-parameter}, differentiate with respect to
$u$.  Direct cancellation gives
\[
 s^{5/2}\frac{\dd}{\dd s}\left(\frac{g(s)}{s^{3/2}}\right)
 =\frac32 k(k-1)u^2\geq0.
\]
The profile and the normalized ratio are continuous at every branch
endpoint.  This proves the claim globally.  For $k=1$, both sides vanish on
the triangle-free interval.  Since $g\geq0$ and $s^{3/2}$ is increasing, the
lemma also shows that $g$ itself is nondecreasing.
\end{proof}


%%%%%%%%%%%%%%%%%%%%%%%%%%%%%%%%%%%%%%%%%%%%%%%%%%%%%%%%%%%%%%%%%%%%%%%%%%%%%%%
\section{A sharp paw lower bound}
%%%%%%%%%%%%%%%%%%%%%%%%%%%%%%%%%%%%%%%%%%%%%%%%%%%%%%%%%%%%%%%%%%%%%%%%%%%%%%%

Let
\[
 \tau(x)=\int_{\Omega^2}
 W(x,y)W(x,z)W(y,z)\dd\mu(y)\dd\mu(z),
 \qquad G=\int_\Omega d(x)\tau(x)\dd\mu(x).
\]
Thus $G$ is the paw homomorphism density.

\begin{lemma}[Fractionally degree-weighted edge]
\label{lem:fractional-edge}
For $0<\alpha\leq1$,
\[
 N_\alpha:=\int_{\Omega^2}W(x,y)d(x)^\alpha d(y)^\alpha\dd\mu^2
 \geq p^{1+2\alpha}.
\]
\end{lemma}

\begin{proof}
Put $\theta=1/(1+2\alpha)$ and
$\eta=\alpha/(1+2\alpha)$.  On the support of $W$,
\[
 W=(Wd(x)^\alpha d(y)^\alpha)^\theta
    (W/d(x))^\eta(W/d(y))^\eta.
\]
H\"older's inequality applies after defining the quotients to be zero on
zero-degree fibres.  Each quotient integral is at most one, and hence
$p\leq N_\alpha^\theta$.
\end{proof}

\begin{lemma}[Sharp paw bias]
\label{lem:sharp-paw}
Every graphon of edge density $p\geq1/2$ satisfies
\begin{equation}
 G\geq p\,g(p).
 \label{eq:sharp-paw}
\end{equation}
\end{lemma}

\begin{proof}
Average the three equal representations obtained by putting the paw leaf at
each triangle vertex.  Arithmetic--geometric mean gives
\[
 G\geq
 \int_{\Omega^3}W(x,y)W(x,z)W(y,z)
       (d(x)d(y)d(z))^{1/3}\dd\mu^3.
\]
Put $M=\int d^{1/3}$, $z=M^3$, and use the probability measure
$\dd\nu=d^{1/3}\dd\mu/M$.  The right side is
\[
 z\,t(K_3,W;\nu).
\]
The edge density $s$ under $\nu$ obeys, by
Lemma~\ref{lem:fractional-edge} with $\alpha=1/3$,
\[
 s\geq s_0:=p\left(\frac pz\right)^{2/3}.
\]
Concavity of $d^{1/3}$ gives $z\leq p$, while $s\leq1$ implies
$s_0\leq1$; moreover $z\leq p$ gives $s_0\geq p$.  The sharp triangle
theorem, monotonicity of $g$, and
Lemma~\ref{lem:normalized-triangle} yield
\begin{align*}
 G&\geq z g(s_0)
 =p^{5/2}\frac{g(s_0)}{s_0^{3/2}}\\
 &\geq p^{5/2}\frac{g(p)}{p^{3/2}}
 =p g(p).
\end{align*}
\end{proof}


%%%%%%%%%%%%%%%%%%%%%%%%%%%%%%%%%%%%%%%%%%%%%%%%%%%%%%%%%%%%%%%%%%%%%%%%%%%%%%%
\section{The low-density interval}
%%%%%%%%%%%%%%%%%%%%%%%%%%%%%%%%%%%%%%%%%%%%%%%%%%%%%%%%%%%%%%%%%%%%%%%%%%%%%%%

For $p\in[1/2,2/3]$, write the sharp triangle profile using part masses
$( (1-y)/2,(1-y)/2,y )$:
\begin{equation}
 p=\frac{1+2y-3y^2}{2},\qquad
 g(p)=\frac32y(1-y)^2.
 \label{eq:low-parameter}
\end{equation}
The relation in \eqref{eq:low-parameter} is increasing for
$0\leq y\leq1/3$, and its value at $y=1/8$ is $77/128>3/5$.
Consequently $p\leq3/5$ gives $0\leq y\leq1/8$.

\begin{lemma}[Low-interval scalar comparison]
\label{lem:low-scalar}
For $p\in[1/2,3/5]$,
\[
 p\,g(p)^2\geq c^2f.
\]
\end{lemma}

\begin{proof}
Substitution from \eqref{eq:low-parameter} gives
\[
 p g(p)^2-c^2f=\frac{y^2}{8}R(y),
\]
where
\[
 R(y)=1+6y-159y^2+756y^3-1521y^4+1422y^5-513y^6.
\]
After $y=u/8$, its degree-six Bernstein coefficients on $0\leq u\leq1$
are
\[
 1,\ \frac98,\ \frac{347}{320},\ \frac{2437}{2560},
 \frac{15909}{20480},\ \frac{19349}{32768},
 \frac{108079}{262144}.
\]
They are all positive.
Since the Bernstein basis is nonnegative on $[0,1]$, this proves
$R(y)>0$ throughout the required interval.
\end{proof}

For fixed $x$, put $F_x(a)=K(x,a)$.  Since $S$ is the kernel of
$T_W^2$, equation \eqref{eq:density-kernels} gives
\begin{equation}
 t(F_{148},W)=\int_\Omega\lVert T_WF_x\rVert_2^2\dd\mu(x).
 \label{eq:hilbert-density}
\end{equation}
Projection onto the constant function, followed by Jensen, gives
\begin{equation}
 t(F_{148},W)
 \geq\int_\Omega
 \left(\int_\Omega d(a)K(x,a)\dd\mu(a)\right)^2\dd\mu(x)
 \geq G^2.
 \label{eq:paw-square}
\end{equation}
Lemma~\ref{lem:sharp-paw} and Lemma~\ref{lem:low-scalar} therefore prove
\[
 t(F_{148},W)\geq p^2g(p)^2\geq pc^2f
\]
throughout $p\in[1/2,3/5]$.


%%%%%%%%%%%%%%%%%%%%%%%%%%%%%%%%%%%%%%%%%%%%%%%%%%%%%%%%%%%%%%%%%%%%%%%%%%%%%%%
\section{The two-dimensional projection}
%%%%%%%%%%%%%%%%%%%%%%%%%%%%%%%%%%%%%%%%%%%%%%%%%%%%%%%%%%%%%%%%%%%%%%%%%%%%%%%

For each $x$, put
\begin{align*}
 s(x)&=\int W(x,a)^2\dd\mu(a),
 &v(x)&=s(x)-d(x)^2,\\
 g(x)&=\int d(a)K(x,a)\dd\mu(a),
 &h(x)&=\int W(x,a)S(x,a)^2\dd\mu(a).
\end{align*}
The functions $\one$ and $W(x,\cdot)-d(x)\one$ are orthogonal, with
squared norms $1$ and $v(x)$.  Their inner products with $T_WF_x$ are
$g(x)$ and $h(x)-d(x)g(x)$.  Bessel's inequality gives
\begin{equation}
 \lVert T_WF_x\rVert_2^2
 \geq g(x)^2+\frac{(h(x)-d(x)g(x))^2}{v(x)},
 \label{eq:bessel}
\end{equation}
where the quotient is zero when $v(x)=0$; its numerator then also vanishes.

Set
\begin{align*}
 D&=\int_{\Omega^2}W(x,y)S(x,y)^2\dd\mu^2,\\
 L&=\int_{\Omega^2}W(x,y)S(x,y)d(x)d(y)\dd\mu^2,\\
 \Delta&=D-L,
 &V&=\int_\Omega v(x)\dd\mu(x).
\end{align*}
Then $\int g=G$ and $\int(h-dg)=\Delta$.  Two applications of
Cauchy--Schwarz in \eqref{eq:bessel} give
\begin{equation}
t(F_{148},W)\geq G^2+\frac{\Delta^2}{V}.
\label{eq:global-projection}
\end{equation}
Here and below the quotient is defined to be zero when $V=0$; in that
case $\Delta=0$ by the preceding fibrewise observation.
Moreover
\[
 V=\int W^2-\int d^2\leq p-p^2=pq.
\]
Thus, for $p<1$,
\begin{equation}
 t(F_{148},W)\geq G^2+\frac{\Delta^2}{pq}.
 \label{eq:coarse-projection}
\end{equation}


%%%%%%%%%%%%%%%%%%%%%%%%%%%%%%%%%%%%%%%%%%%%%%%%%%%%%%%%%%%%%%%%%%%%%%%%%%%%%%%
\section{The high-density supporting line}
%%%%%%%%%%%%%%%%%%%%%%%%%%%%%%%%%%%%%%%%%%%%%%%%%%%%%%%%%%%%%%%%%%%%%%%%%%%%%%%

\begin{lemma}[High-density linear estimate]
\label{lem:linear}
If $p\in[3/5,1]$, then
\begin{equation}
 p^2G-q\Delta\geq pcf.
 \label{eq:linear}
\end{equation}
\end{lemma}

\begin{proof}
Symmetry and arithmetic--geometric mean give
\[
 G=\int WS\frac{d(x)+d(y)}2\geq\int WSZ,
 \qquad Z=\sqrt{d(x)d(y)}.
\]
Also $L=\int WSZ^2$ and $D=\int WS^2$.  Hence the left side of
\eqref{eq:linear} is at least
\begin{equation}
 \int_{\Omega^2}W(x,y)F_p(Z,S)\dd\mu^2,
 \qquad
 F_p(z,s)=s(p^2z+qz^2-qs).
 \label{eq:F}
\end{equation}
Cauchy--Schwarz and the product union bound give
\[
 (2z-1)_+\leq S\leq z.
\]
For fixed $z$, $F_p(z,s)$ is concave in $s$, so it is bounded below by
the smaller of its endpoint values
\begin{align*}
 A_p(z)&=(2z-1)_+
 \bigl(p^2z+qz^2-q(2z-1)_+\bigr),\\
 B_p(z)&=z^2(p^2-q+qz).
\end{align*}

For $z\geq1/2$, put
\[
 \ell_p(z)=cf+m_p(z-p),
 \qquad m_p=-2p^3+15p^2-14p+4.
\]
On $p\in[3/5,1]$, one has $m_p\geq0$: after substituting
$p=3/5+2a/5$, its degree-three Bernstein coefficients on
$0\leq a\leq1$ are
\[
 \frac{71}{125},\quad \frac{61}{75},\quad \frac53,\quad 3.
\]
Also
\[
 \ell_p(1/2)=\frac{c^2}{2}(p^2-4p+2)\leq0.
\]
Furthermore
\[
 A_p(z)-\ell_p(z)
 =-(z-p)^2\bigl(2p^2+2pz-9p-2z+5\bigr)\geq0,
\]
because the bracket is decreasing in $z$ and at $z=1/2$ equals
$2(p^2-4p+2)\leq0$.  Indeed, $p^2-4p+2$ is decreasing on this
interval and has value $-1/25$ at $p=3/5$.

The remaining polynomial $B_p(z)-\ell_p(z)$ is nonnegative on the full
rectangle $3/5\leq p\leq1$, $0\leq z\leq1$.  For completeness, substitute
$p=3/5+2a/5$ and affinely map each of the four $z$-intervals
\[
 [0,1/2],\quad[1/2,2/3],\quad[2/3,3/4],\quad[3/4,1]
\]
to $0\leq b\leq1$.  Its native bidegree is $(4,3)$, and its four exact
Bernstein coefficient matrices are
\[
\begin{pmatrix}
178/625&713/3750&691/7500&51/1250\\
48/125&97/375&89/600&181/2000\\
2/3&23/50&13/45&107/600\\
6/5&13/15&71/120&31/80\\
2&3/2&13/12&3/4
\end{pmatrix},
\]
\[
\begin{pmatrix}
51/1250&533/22500&397/33750&116/16875\\
181/2000&641/9000&13/225&58/1125\\
107/600&191/1350&451/4050&4/45\\
31/80&23/72&7/27&28/135\\
3/4&23/36&29/54&4/9
\end{pmatrix},
\]
\[
\begin{pmatrix}
116/16875&299/67500&337/90000&101/20000\\
58/1125&109/2250&1697/36000&153/3200\\
4/45&629/8100&41/600&293/4800\\
28/135&49/270&227/1440&87/640\\
4/9&43/108&17/48&5/16
\end{pmatrix},
\]
\[
\begin{pmatrix}
101/20000&269/30000&77/2500&48/625\\
153/3200&299/6000&209/3000&14/125\\
293/4800&47/1200&8/225&4/75\\
87/640&17/240&1/40&0\\
5/16&3/16&1/12&0
\end{pmatrix}.
\]
Every entry is nonnegative.
The tensor-product Bernstein basis is nonnegative and sums to one on each
box, so these matrices prove $B_p(z)\geq\ell_p(z)$ on the full rectangle.
Thus $F_p(z,s)\geq\ell_p(z)$ also when
$z<1/2$, since $A_p(z)=0\geq\ell_p(z)$ there.

Finally use the edge-biased probability measure
$W(x,y)\dd\mu^2/p$.  The edge geometric-mean inequality gives
\[
 \mathbb E\sqrt{d(x)d(y)}\geq p.
\]
Indeed, after defining zero-degree quotients to be zero,
Cauchy--Schwarz applied to
$\sqrt{W/d(x)}$ and $\sqrt{W/d(y)}$ gives
$\int W/Z\leq1$.  A second application of Cauchy--Schwarz gives
$p^2\leq(\int WZ)(\int W/Z)$.  Since $m_p\geq0$, integration of the
supporting line proves
\[
 \int WF_p(Z,S)\geq p\ell_p(\mathbb E Z)\geq pcf.
\]
\end{proof}

\begin{lemma}[Projection closes the high interval]
For $p\in[3/5,1]$,
\[
 t(F_{148},W)\geq pc^2f.
\]
\end{lemma}

\begin{proof}
The case $p=1$ is immediate.  For $p<1$, put
\[
 G_0=p^2c,\qquad \Delta_0=-pcq^2.
\]
Then
\[
 G_0^2+\frac{\Delta_0^2}{pq}=pc^2(p^3+q^3)=pc^2f.
\]
Lemma~\ref{lem:linear} is precisely
\[
 GG_0+\frac{\Delta\Delta_0}{pq}\geq pc^2f.
\]
Cauchy--Schwarz in this weighted two-dimensional Euclidean space and
\eqref{eq:coarse-projection} finish the proof.
\end{proof}


%%%%%%%%%%%%%%%%%%%%%%%%%%%%%%%%%%%%%%%%%%%%%%%%%%%%%%%%%%%%%%%%%%%%%%%%%%%%%%%
\section{Conclusion and scope}
%%%%%%%%%%%%%%%%%%%%%%%%%%%%%%%%%%%%%%%%%%%%%%%%%%%%%%%%%%%%%%%%%%%%%%%%%%%%%%%

\begin{theorem}[Atlas 148, graph6 \texttt{EyUG}]
Every graphon $W$ of edge density $p\in[1/2,1]$ satisfies
\[
 \boxed{
 t(F_{148},W)\geq p(2p-1)^2(3p^2-3p+1)
 =\Phi_{F_{148}}(p).
 }
\]
\end{theorem}

The low- and high-density arguments meet at $p=3/5$.  At $p=1/2$ the
target is zero (with equality on the balanced complete bipartite graphon),
and at $p=1$ both sides are one.  No classification of all interior equality
cases is asserted.  All integral identities use
bounded nonnegative measurable functions, except for ordinary Hilbert-space
inner products of bounded kernels, so Tonelli--Fubini applies on an arbitrary
probability space.  Zero-degree and zero-variance fibres were handled before
division.

This proof does not use the false comparison
$p^2t(F_{148})\geq t(K_3)^2t(C_4)$.  The two-coordinate projection retains
exactly the covariance term that comparison discards.  The use of the sharp
triangle-density theorem is confined to the paw estimate on
$[1/2,3/5]$; the high interval uses arithmetic--geometric mean,
Cauchy--Schwarz, and the displayed rational Bernstein certificate.

The Atlas identification, chromatic target, normalized-profile derivative,
low-interval factorization, projection identities, and all Bernstein
coefficients are reconstructed exactly by
\begin{verbatim}
python codes/verify_atlas148_paw_bias_hilbert_projection.py
\end{verbatim}

No Lean file is changed.  Any future formalization would need either the
graphon triangle-density theorem as an imported analytic input or a separate
formal proof of its $K_3$ case, as well as the two-dimensional Bessel and
Bernstein steps above.

\begin{thebibliography}{1}
\bibitem{ReiherCliqueDensity}
C.~Reiher,
\emph{The clique density theorem},
Annals of Mathematics \textbf{184} (2016), 683--707;
\url{https://annals.math.princeton.edu/2016/184-3/p01}.
\end{thebibliography}


%%%%%%%%%%%%%%%%%%%%%%%%%%%%%%%%%%%%%%%%%%%%%%%%%%%%%%%%%%%%%%%%%%%%%%%%%%%%%%%
\appendix
\section{What the Lean formalization actually proves}
%%%%%%%%%%%%%%%%%%%%%%%%%%%%%%%%%%%%%%%%%%%%%%%%%%%%%%%%%%%%%%%%%%%%%%%%%%%%%%%

The theorem above is formalized in full, on the whole interval $p\in[1/2,1]$.
The development is \texttt{lean/Taeyoung/Methods/Atlas148/}, the catalogue
theorem is \texttt{satisfiesLowerBound\_148}, and every declaration depends on
exactly \texttt{[propext, Classical.choice, Quot.sound]}.  This supersedes the
closing remark of Section~7: no Lean file was left unchanged, and the two
steps flagged there as obstacles turned out not to be needed in the stated
form.  This appendix says where the formalization departs from the argument
above.

\subsection{Only Fisher's band of the clique density theorem is used}
\label{sec:app148-fisher}

Section~2 imports the full clique density theorem, and the paw estimate applies
it at the tilted edge density $s$, which ranges over all of $[p,1]$.  Only the
$k=2$ branch is needed.  Split at $s=2/3$:

\begin{itemize}
 \item for $s\leq2/3$ the sharp profile is Fisher's 1989 theorem, whose
 parameter $c=(1-\sqrt{4-6s})/3$ is exactly the $y$ of
 \eqref{eq:low-parameter};
 \item for $s\geq2/3$ the crude Goodman bound $t(K_3)\geq s(2s-1)$ already
 suffices.  Fed through the same substitution $z=p^{5/2}s^{-3/2}$ it gives
 $p^{5/2}s^{-1/2}(2s-1)$, increasing in $s$, so its minimum over $[2/3,1]$ is
 at $s=2/3$; and there it clears the requirement $p\,g(p)$ by a factor of at
 least $1.34$ on $[1/2,3/5]$.
\end{itemize}
The two agree at $s=2/3$, where both equal $2/9$, so no case is lost.  The
Goodman branch reduces to $6g(p)^2\leq p^3$, i.e.\ $108y^2(1-y)\leq(1+3y)^3$,
whose difference factors as $(3y-1)^2(15y+1)$.

Fisher's theorem is nevertheless unavoidable: the triangle input must be
within $3.5\%$--$6\%$ of sharp, the required ratio $c\sqrt{f/p}\,/\,g(p)$
peaking at $0.9656$ at $p=3/5$ and equalling $0.9428$ at $p=1/2$, whereas
Goodman alone runs $0.667$--$0.848$ of sharp.  The Lean proof imports it from
\texttt{lean/Taeyoung/Fisher/}, a copy of the sibling project
\texttt{discussions/goodman-style-bound/fisher\_lean} restricted to the closure
of its target; \texttt{Methods/TriangleDensity.lean} restates it for this
project's bundled graphons.

\subsection{The fractional edge lemma needs no three-factor H\"older}
\label{sec:app148-holder}

Lemma~\ref{lem:fractional-edge} is proved above by H\"older at three factors,
with the unbounded quotients $W/d(x)$.  One \emph{two}-factor Cauchy--Schwarz
does it.  With $u=d^{1/3}$,
\[
 p^2=\Bigl(\int\!\!\int W\Bigr)^2
 \leq\Bigl(\int\!\!\int \frac{W}{u_xu_y}\Bigr)
      \Bigl(\int\!\!\int W u_xu_y\Bigr)
 \leq M\cdot N_{1/3},
\]
the first factor being at most $M=\int d^{1/3}$ by arithmetic--geometric mean
and Fubini.  Combined with $M\leq p^{1/3}$ — Jensen for the concave cube root
— this gives $p^2\leq p^{1/3}N_{1/3}$, hence $p^5\leq N_{1/3}^3$.

Unboundedness is handled by shifting the degree to $d+\delta^3$, so that
$(d+\delta^3)^{1/3}\leq d^{1/3}+\delta$, and letting $\delta\to0$; the same
device gives the edge geometric mean $\int\!\!\int W\sqrt{d(x)d(y)}\geq p^2$
of Section~6, which the note proves through the unbounded $\int\!\!\int W/Z$.

Everything downstream is kept polynomial: the tilt sends $z=M^3$ and
$s=N_{1/3}/M^2$, so the hypothesis $p^5\leq z^2s^3$ that the scalar lemmas
consume is \emph{identically} $p^5\leq N_{1/3}^3$, the $M$'s cancelling.

\subsection{Bessel without division, and the final Cauchy--Schwarz for free}
\label{sec:app148-bessel}

The quotient form \eqref{eq:bessel} is never used.  Keeping the multiplier
free, for every real $\lambda$,
\[
 \lVert T_WF_x\rVert^2\geq g(x)^2+2\lambda\bigl(h(x)-d(x)g(x)\bigr)
 -\lambda^2v(x),
\]
which needs neither a division nor the separate treatment of the degenerate
fibres $v(x)=0$.  In that form Bessel is not a Hilbert-space statement at all:
putting $w=T_WF_x-\lambda(W(x,\cdot)-d(x))$, the centred function has mean
zero, so $\int w=g(x)$, and plain Jensen $(\int w)^2\leq\int w^2$ is the whole
content.

The same free $\lambda$ removes the two-dimensional Cauchy--Schwarz of
Section~7.  Integrating gives $T\geq G^2+2\lambda\Delta-\lambda^2V$ with
$V\leq pq$; taking $\lambda=-cq$ and inserting the linear estimate
\eqref{eq:linear} turns the target into
\[
 T-pc^2f\geq(G-cp^2)^2\geq0,
\]
because $f-q^3=p^3$.  So the vector $(G_0,\Delta_0)$, the weighted inner
product, and the division by $V$ all disappear; $V\leq pq$ survives only as a
coefficient bound.

\subsection{Less Bernstein data, and no permuted triple integral}
\label{sec:app148-bernstein}

The note certifies three scalar facts with $91$ Bernstein coefficients: seven
for Lemma~\ref{lem:low-scalar}, four for the slope $m_p\geq0$, and eighty for
$B_p\geq\ell_p$ on four $z$-boxes.  The formalization uses $40$.

\begin{itemize}
 \item Lemma~\ref{lem:low-scalar} and the normalized monotonicity of
 Lemma~\ref{lem:normalized-triangle} are closed directly by \texttt{nlinarith}
 in the coordinate $y$, the latter after cross-multiplying and squaring into a
 two-variable polynomial inequality on $0\leq y_1\leq y_2\leq1/3$.  The
 normalized-profile derivative also collapses: on the $k=2$ branch,
 $g(s)^2/s^3=18y^2(1-y)/(1+3y)^3$, with derivative
 $36y(1-3y)/(1+3y)^4$.
 \item For $B_p\geq\ell_p$ the subdivision is genuinely necessary — a single
 box fails at every bidegree up to $(10,12)$, because the Bernstein
 coefficients converge only like $O(1/n)$ along the edge $p=3/5$, where the
 function's minimum is a shallow $0.0044$; and a linear-programming search over
 all products of the four box constraints, with $(z-p)^2$ multiples adjoined,
 is infeasible.  But \emph{two} boxes suffice, split at $z=7/10$ and kept at
 the native bidegree $(4,3)$.  Each is discharged by \texttt{linarith} against
 the twenty products $a^i(1-a)^{4-i}b^j(1-b)^{3-j}$, so no certificate
 datatype is involved.
\end{itemize}

The leaf symmetrization of Lemma~\ref{lem:sharp-paw} needs no permutation of a
triple integral.  Integrating the third vertex first turns each of the three
degree weights into a page operator
$H_s(x,y)=\int W(x,z)W(y,z)d(z)^s$, so arithmetic--geometric mean becomes the
pointwise estimate
\[
 d(x)^{1/3}d(y)^{1/3}H_{1/3}(x,y)
 \leq\frac{d(x)+d(y)}{3}H_0(x,y)+\frac{H_1(x,y)}{3},
\]
and one monotone step over the spine pair finishes it; the three terms on the
right are all $G$, two by symmetry of $W\!\cdot\!H_0$ and the third because
$\int\!\!\int W H_1$ and $G$ both collapse to $\int d\,\tau$.

\subsection{Two things that are load-bearing}
\label{sec:app148-sharp}

\begin{itemize}
 \item \textbf{The collapse to $Z=\sqrt{d(x)d(y)}$.}  Keeping $d(x)$ and
 $d(y)$ separate and replacing $Z$ by $(d(x)+d(y))/2$ in the supporting line
 makes it \emph{false}: it is violated by $-0.014$ at $p=3/5$, $d(x)=0.28$,
 $d(y)=1$.  What forces $Z$ is that $L=\int\!\!\int W S Z^2$ holds exactly,
 whereas $d(x)d(y)\leq\bar d^2$ points the wrong way.
 \item \textbf{The split point.}  The high-density argument needs
 $\ell_p(1/2)\leq0$, that is $p^2-4p+2\leq0$, so it cannot reach below
 $p=2-\sqrt2\approx0.5858$.  The window $[1/2,2-\sqrt2]$ is irreducibly
 Fisher-dependent; a supporting line valid down to $p=1/2$ would remove the
 triangle-density input from the catalogue entirely.
\end{itemize}

\subsection{What is checked}
\label{sec:app148-scope}

The Lean graph is labelled for peeling — centre $0$, the two joined outer
vertices $1$ and $2$, the triangle apexes $3$ and $4$, the path middle $5$ —
and \texttt{Examples/Graph148.lean} transports the bound to the Atlas edge list
along an explicit relabelling whose adjacency correspondence is checked by
kernel computation.  The chromatic polynomial \eqref{eq:chromatic} is proved,
not assumed: its last factor $r^2-3r+3$ is irreducible over the rationals, so
no clique-attachment tower produces it, and the six surjective colouring counts
$0,0,0,18,264,840,720$ are checked by kernel reduction instead, the top one
through the identity $\mathrm{surj}(n)=n!$.

\end{document}
